\section*{Chapter \ref{ch:b2:fragment-orbitals} --- Fragment Orbitals and Polyatomic Molecules}

\begin{solution}{exo:b2:fragment-orbitals:1}
$h_1 + h_2$ and $h_1 - h_2$. The plane $xz$, which bisects the angle and is
perpendicular to the molecule, exchanges the two hydrogens: $h_1 - h_2$ changes
sign, $h_1 + h_2$ does not.
\end{solution}

\begin{solution}{exo:b2:fragment-orbitals:2}
Six valence orbitals (oxygen 2s and three 2p, two hydrogen 1s), eight valence
electrons, four filled orbitals.
\end{solution}

\begin{solution}{exo:b2:fragment-orbitals:3}
The filled valence levels are one $a_1$ orbital and a set of three degenerate $t_2$
orbitals: two energies, two bands. The $t_2$ band removes electrons from three
orbitals holding six electrons, against one orbital and two electrons for $a_1$.
\end{solution}

\begin{solution}{exo:b2:fragment-orbitals:4}
The $\pi$ bond needs the two p orbitals of the carbons to be parallel, which places
both \ce{CH2} planes in a common plane.
\end{solution}

\begin{solution}{exo:b2:fragment-orbitals:5}
Ammonia loses an electron from its lone pair, a mostly \omterm{def:b2:diatomic-mos:bonding}{nonbonding orbital} on
nitrogen; methane from a bonding $t_2$ orbital. A nonbonding electron lies higher,
closer to the atomic level, than a bonding one: it is removed more easily.
\end{solution}

\begin{solution}{exo:b2:fragment-orbitals:6}
Linear water has two degenerate filled nonbonding p orbitals; bending lowers one of
them and slightly raises a bonding level: a net gain with eight electrons. \ce{BeH2}
has only four, which fill the two \omterm{def:b2:diatomic-mos:bonding}{bonding orbitals}; the orbital that bending would
lower is empty, while the bonding levels lose overlap: the linear shape is best.
\end{solution}

\begin{solution}{exo:b2:fragment-orbitals:7}
Each C--H or C--C bond on a carbon next to the cationic centre can align with the
empty p orbital and gives about $2\beta^2/\Delta$ of stabilisation. \ce{CH3+} has no such
neighbour; ethyl, isopropyl and tert-butyl cations have one, two and three methyl
groups next to the centre, each offering an aligned bond: stability grows in that
order.
\end{solution}

\begin{solution}{exo:b2:fragment-orbitals:8}
For two equal levels the splitting is $2|\beta|$ and the two $\pi$ electrons gain $2|\beta|$,
proportional to $\cos\theta$. It falls to half at $\theta = 60^\circ$ and to zero at
$90^\circ$.
\end{solution}

\begin{solution}{exo:b2:fragment-orbitals:9}
Borane has six valence electrons: they fill the three \omterm{def:b2:diatomic-mos:bonding}{bonding orbitals}, and the p
orbital of boron perpendicular to the plane stays empty; nothing is gained by
pyramidalising. Ammonia has two more: they occupy the orbital that is nonbonding in
the planar shape, and, as in water, pyramidalising lets it mix with the 2s orbital and
the hydrogen combination and drop.
\end{solution}

\begin{solution}{exo:b2:fragment-orbitals:10}
The secular equations for $c_1 = c_2 = c_3$ give $\alpha + 2\beta$; the two combinations
orthogonal to it give $\alpha - \beta$ (degenerate). Two electrons in $\alpha + 2\beta$: energy
$2\alpha + 4\beta$, lower than $\ce{H2} + \ce{H+}$ ($2\alpha + 2\beta$) since $\beta < 0$.
\end{solution}

\begin{solution}{exo:b2:fragment-orbitals:11}
$\pi_1$: all three coefficients of one sign, no node (bonding); $\pi_2$: opposite signs
on the end carbons and a node at the centre (nonbonding); $\pi_3$: alternating
signs, two nodes (antibonding). The anion has four $\pi$ electrons: $\pi_2$ is its
highest occupied level, with density only on the end carbons.
\end{solution}

\begin{solution}{exo:b2:fragment-orbitals:12}
In each of the two planes containing the axis, three parallel p orbitals (O, C, O)
give a bonding, a nonbonding (on the oxygens only) and an \omterm{def:b2:diatomic-mos:bonding}{antibonding orbital}: six
$\pi$ orbitals in all. The two bonding and the two nonbonding ones are filled: eight
$\pi$ electrons (two $\pi$ bonds and two oxygen lone pairs in the Lewis picture).
\end{solution}

\begin{solution}{pb:b2:fragment-orbitals:1}
\textbf{1.} Eight: six from oxygen, one from each hydrogen.
\textbf{2.} $h_1 + h_2$ and $h_1 - h_2$.
\textbf{3.} 2s and $2\mathrm p_z$.
\textbf{4.} $2\mathrm p_y$.
\textbf{5.} $2\mathrm p_x$, perpendicular to the molecular plane.
\textbf{6.} Six; four filled.
\textbf{7.} $h_1 - h_2$, which changes sign, like $2\mathrm p_z$, through the plane
perpendicular to the axis at oxygen.
\textbf{8.} $h_1 + h_2$.
\textbf{9.} $2\mathrm p_x$ and $2\mathrm p_y$, perpendicular to the axis: a degenerate
nonbonding pair.
\textbf{10.} (2s, $h_1 + h_2$) bonding$^2$, ($2\mathrm p_z$, $h_1 - h_2$) bonding$^2$, then
$2\mathrm p_x^2\,2\mathrm p_y^2$.
\textbf{11.} No: the degenerate pair is full, every electron is paired.
\textbf{12.} At the energy of an oxygen 2p orbital: it is nonbonding.
\textbf{13.} $2\mathrm p_y$ (the axis of the linear molecule along $z$, bending in $yz$).
\textbf{14.} It mixes with 2s and $h_1 + h_2$ and drops: the $3a_1$ level.
\textbf{15.} The ($2\mathrm p_z$, $h_1 - h_2$) bonding level, which loses some overlap as the
hydrogens leave the axis: it becomes $1b_2$.
\textbf{16.} The filled level that drops gains more than the other loses: with eight
electrons the bent shape is the more stable.
\textbf{17.} $104.48^\circ$, between $90^\circ$ (bonds made of p orbitals only) and $109.5^\circ$:
the 2s orbital takes some part in the bonding.
\textbf{18.} $2\mathrm p_x$, which stays the nonbonding $1b_1$.
\textbf{19.} Four.
\textbf{20.} From $1b_1$, the nonbonding $2\mathrm p_x$ orbital of oxygen.
\textbf{21.} Removing a nonbonding electron hardly changes the bonds: the ion has the
geometry of the molecule and little vibrational energy is excited. Removing a
bonding electron changes the bond lengths and angles, and the band spreads over
many vibrational levels.
\textbf{22.} \qty{12.62}{eV} against \qty{13.62}{eV}: oxygen in water carries a partial
negative charge drawn from the hydrogens, which raises its 2p levels.
\textbf{23.} They are not two equal levels: their two combinations are the $3a_1$ and
$1b_1$ orbitals, of different energies. Both pictures describe the same total
electron density; only the delocalised orbitals give the energies of ionisation.
\textbf{24.} Water has \textbf{four} valence ionisation bands, the lowest at
$\boldsymbol{\approx \qty{12.6}{eV}}$.
\end{solution}
